\documentclass[11pt]{article}
\usepackage[T1]{fontenc}
\usepackage{lmodern,amsmath,amssymb,amsthm,booktabs,microtype}
\usepackage[margin=1in]{geometry}
\usepackage[hidelinks]{hyperref}
\hypersetup{pdftitle={Toward a Universal Model: Preparation, matched measurements, and relative phase}}
\newtheorem{proposition}{Proposition}[section]
\newtheorem{corollary}[proposition]{Corollary}
\DeclareMathOperator{\Ran}{Ran}
\DeclareMathOperator{\Span}{span}
\DeclareMathOperator{\rank}{rank}
\DeclareMathOperator{\Tr}{Tr}
\newcommand{\C}{\mathbb C}
\newcommand{\R}{\mathbb R}
\newcommand{\Hmin}{\mathcal H_{\min}}
\newcommand{\ket}[1]{\lvert #1\rangle}
\begin{document}
\begin{center}
{\LARGE Toward a Universal Model\par}
\smallskip
{\large Preparation, matched measurements, and relative phase\par}
\smallskip
Working manuscript
\end{center}
\begin{abstract}
\noindent
A preparation map selects a subspace of an output Hilbert space. We characterize the further directions generated by evolution of the prepared states and count them from a finite positive response. A four-mode fermion model exhibits complete transfer to one such direction. An exact perturbation calculation shows how an occupation measurement can record near-complete transfer even when appreciable weight has leaked outside the selected carrier. On the directly prepared space the polar isometry carries input states and measurement effects to pairs with equal Born probabilities. The same occupation detector illustrates both this correspondence and its dynamical scope. On the direct sum of the active supports we characterize the measurement effects that distinguish relative phases and show how evolution can preserve, rotate, or leave the associated graphs.
\end{abstract}

\section{Preparation, evolution, and the final defect}
\label{sec:preparation-evolution}
\subsection{The directly prepared space}
Let $A:\operatorname{Dom}A\subseteq\mathcal H_\Psi\to\mathcal H_\Phi$ be a densely defined closed linear map between Hilbert spaces. In the polar decomposition $A=u|A|$, the positive operator $|A|$ weights input directions and the partial isometry $u$ transports directions on $(\ker A)^\perp$. The directly prepared space and the final defect are
\[
 D=\overline{\Ran A}=\Ran u,\qquad D^\perp=\ker A^*=\ker u^* .
\]
Enlarging the codomain adds its orthogonal complement to the final defect. Deciding which part of the defect is accessible requires dynamics.

Let $H$ be self-adjoint in the positive Hilbert inner product on $\mathcal H_\Phi$. Write $U(t)=\exp(-itH/\hbar)$ and set
\[
 \Hmin=\overline{\Span}\{U(t)v:v\in D,\ t\in\R\},\qquad
 \mathcal D=\Hmin\ominus D .
\]
The closed span permits linear combinations and limits. Complete transfer to a prescribed ray in a single wait is a further property of the model.

\begin{proposition}[Decomposition for a fixed preparation]\label{prop:preparation-split}
The output space is the orthogonal sum
\[
 \mathcal H_\Phi=D\oplus\mathcal D\oplus\Hmin^\perp.
\]
For any bounded port $\Gamma:\mathcal E\to\mathcal H_\Phi$ with $\overline{\Ran\Gamma}=D$, the vectors invisible to its reciprocal readout at every real time satisfy
\[
 \{y:\Gamma^*U(t)y=0\text{ for every }t\in\R\}=\Hmin^\perp.
\]
The adjoint kernel of $A$ with its full codomain is $\mathcal D\oplus\Hmin^\perp$. With codomain restricted to $\Hmin$, it is $\mathcal D$.
\end{proposition}
\begin{proof}
Taking $t=0$ gives $D\subseteq\Hmin$. The unitary group preserves $\Hmin$ in both time directions. Also $\Gamma^*U(t)y=0$ precisely when $y$ is orthogonal to $U(-t)\Ran\Gamma$. Taking all real $t$ and closing the span proves the assertion. The two kernel statements follow by taking the orthogonal complement of $D$ in the indicated codomain.
\end{proof}
In finite dimensions this is Kalman's duality between controllability and observability applied to a self-adjoint system \cite[Sec.~5]{Kalman1963}.

\subsection{Forward waiting and the measured response}

Forward preparation generates $\mathcal H_{\mathrm{fwd}}=\overline{\Span}\{U(t)v:v\in D,\ t\geq0\}$. Forward detection involves $U(-t)D$ for $t\geq0$, by the adjoint calculation above.

\begin{proposition}[Waiting with a semibounded output Hamiltonian]\label{prop:semibounded-waiting}
If the spectrum of the self-adjoint output Hamiltonian $H$ is bounded below or bounded above, then every nonempty open time interval $\mathcal I$ generates the same closed linear span:
\[
 \overline{\Span}\{U(t)v:v\in D,\ t\in\mathcal I\}=\Hmin .
\]
In particular the forward and backward orbit spans both equal $\Hmin$.
\end{proposition}
\begin{proof}
We use the boundary uniqueness argument of Hegerfeldt \cite[Sec.~2]{Hegerfeldt1998}. Suppose $H\geq aI$ and put $K=H-aI$. This shift multiplies each real-time orbit vector by a scalar phase. Take $y$ orthogonal to the interval span and $v\in D$. With the inner product linear in its second entry, set $F(z)=\langle y,\exp(-izK/\hbar)v\rangle$. The spectral multiplier has modulus at most one on the closed lower half-plane and decays exponentially at high energy in its interior. The spectral theorem then gives boundedness, interior holomorphy and continuity up to the real axis. On $\mathcal I$ the boundary values of $F$ are zero. Schwarz reflection extends $F$ holomorphically across an open segment of $\mathcal I$. The identity theorem gives zero on the lower half-plane, and continuity gives zero at every real time. Thus $y\perp\Hmin$. The reverse inclusion of orthogonal complements follows from the definition. For $H$ bounded above, apply the result to $-H$ on $-\mathcal I$.
\end{proof}
An arbitrarily short interval can therefore generate the full closed span even when the amplitude on a particular direction is small throughout that interval. Transfer probabilities and the waiting time for complete transfer require their own calculation.

For translations on $L^2(\R)$ the generator has spectrum $\R$. Write $(U(t)f)(x)=f(x-t)$, take $D$ to consist of functions supported on $[0,\infty)$, and use its inclusion as the port. A positive wait shifts every preparation to the right. Thus $\mathcal H_{\mathrm{fwd}}=D$. Translations through all real times generate $L^2(\R)$, since the union of the subspaces supported on $[-n,\infty)$ is dense. A nonzero function supported on the negative half-line can enter the port under a positive wait even though forward preparation has zero weight there.

The criterion for $\mathcal D=\{0\}$ is $U(t)D=D$ for every $t\in\R$. For bounded $H$ and an isometric preparation $A$, put $\Pi_D=AA^*$. The equivalent generator condition is
\begin{equation}
 (I-\Pi_D)HA=0,\qquad HA=A(A^*HA).
 \label{eq:finite-preparation-invariance}
\end{equation}
The first equality says $HD\subseteq D$. Self-adjointness then makes $D$ reducing, and the exponential preserves $D$ for all $t$. Conversely, differentiation at zero gives invariance under $H$. The induced input generator is $A^*HA$. Intertwining a specified input Hamiltonian requires that it equal this compression.

For a bounded port, the measured response is the compressed resolvent $M_\Gamma(z)=\Gamma^*(H-zI)^{-1}\Gamma$, $\operatorname{Im}z\ne0$. The polar port $u$ is available even for unbounded $A$. Replacing $u$ by $2u$ preserves all the subspaces in Proposition~\ref{prop:preparation-split} and multiplies the response by four.

\begin{corollary}[Dimension determined by the complete finite response]\label{cor:finite-response-defect}
Let $H$ be self-adjoint on $\mathcal H_\Phi$ and let $\Gamma:\mathcal E\to\mathcal H_\Phi$ be a bounded port with $\overline{\Ran\Gamma}=D$. Suppose its full compressed spectral measure consists of finitely many distinct atoms $\lambda_k$ with positive finite-rank weights $R_k=\Gamma^*E_H(\{\lambda_k\})\Gamma$. Then
\[
 \dim\mathcal D=\sum_k\rank R_k-\rank\!\left(\sum_kR_k\right).
\]
\end{corollary}
\begin{proof}
Write $E_k=E_H(\{\lambda_k\})$. Full spectral mass gives $\sum_k E_k\Gamma=\Gamma$. Each $E_k\Gamma$ has rank $\rank R_k$ since $R_k=(E_k\Gamma)^*E_k\Gamma$. Their ranges are orthogonal and span $\Hmin$: each frequency component is isolated by averaging $\exp(it\lambda_k/\hbar)U(t)\Gamma h$ over $[0,T]$ as $T\to\infty$. Their ranks therefore add to the minimal dimension, Gilbert's simple-pole formula \cite[Thm.~7]{Gilbert1963}. The direct range has dimension $\rank(\Gamma^*\Gamma)=\rank(\sum_kR_k)$. Subtracting proves the count.
\end{proof}

\subsection{A ground and two excitation energies}

First consider an isometry that takes an orthonormal input basis to a ground $g$ and the excitation $b=(f_1+f_2)/\sqrt2$. The $f_i$ are orthonormal energy eigenstates above $g$. If their energies $E_1,E_2$ differ, the ground contributes a rank-one weight and each excited line contributes half the same rank-one input projection. These weights sum to $I$, giving $\dim\mathcal D=3-2=1$. At $t=\pi\hbar/|E_1-E_2|$ the relative phase reverses and $b$ reaches the ray of $(f_1-f_2)/\sqrt2$. Equal energies merge the two excited weights and give $2-2=0$.

A pair Hamiltonian realizes this calculation with a specified occupation measurement. Following the canonical component summarized by Singh \cite[Sec.~3.4, Eqs.~(3.11)--(3.12)]{Singh2026PairTransfer}, assume four canonical fermion modes $c_1,c_2,d_1,d_2$ with their ordinary Hilbert adjoints. With $P_c=c_2c_1$, $P_d=d_2d_1$, write the symmetric family on the six-dimensional two-particle space $\mathcal F_2$ as
\[
 H_s=-s(P_c^*P_c+P_d^*P_d)-j(P_c^*P_d+P_d^*P_c),\qquad j>0,\quad s\geq0 .
\]
The bare component has $s=0$. The exchange coefficient matrix has eigenvalues $s\pm j$, so its minimal positive semidefinite completion in this family is $s=j$. Write $C=P_c^*\ket{\mathrm{vac}}$, $D_p=P_d^*\ket{\mathrm{vac}}$, $g=(C+D_p)/\sqrt2$, $e=(C-D_p)/\sqrt2$, and
\[
 m=(c_1^*d_2^*+d_1^*c_2^*)\ket{\mathrm{vac}}/\sqrt2,\qquad
 W=\mathcal F_2\ominus\Span\{g,e,m\}.
\]
Direct action gives $H_sg=-(s+j)g$, $H_se=(j-s)e$ and $H_s=0$ on $\C m\oplus W$. Thus $g$ is the unique ground.

The apparatus prepares the two pair occupations after the same orbital rotation on each internal label:
\[
 f_i^*=\cos\theta\,c_i^*+\sin\theta\,d_i^*,\quad
 h_i^*=-\sin\theta\,c_i^*+\cos\theta\,d_i^*,\qquad 0<\theta<\pi/4.
\]
The isometry $A_\theta$ sends the input pair basis to $f_1^*f_2^*\ket{\mathrm{vac}}$ and $h_1^*h_2^*\ket{\mathrm{vac}}$, respectively. Let $g_\Psi,e_\Psi$ be the normalized sum and difference of those input basis states. Expanding the two preparations, with $u_\theta=\cos(2\theta)$ and $v_\theta=\sin(2\theta)$, gives
\[
 A_\theta g_\Psi=g,\quad A_\theta e_\Psi=b=u_\theta e+v_\theta m,\quad
 d=-v_\theta e+u_\theta m .
\]
For the input generator given by restricting $H_s$ to the original pair plane, $g_\Psi$ is its ground. The equalities $A_\theta g_\Psi=g$ and $A_\theta^*g=g_\Psi$ give the ground correspondence. The rotation and its angle specify the preparation interface.

For $s\ne j$, the three port weights at energies $-(s+j), j-s,0$ are $\Pi_{g_\Psi},u_\theta^2\Pi_{e_\Psi},v_\theta^2\Pi_{e_\Psi}$, with $\Pi_x$ the projection onto a unit vector $x$. Their sum is $I$, and the excitation weights are equal at $\theta=\pi/8$. Projecting $H_sb=(j-s)u_\theta e$ orthogonally to the direct plane gives
\begin{equation}
 (I-A_\theta A_\theta^*)H_sA_\theta e_\Psi=-(j-s)u_\theta v_\theta d,\qquad
 (I-A_\theta A_\theta^*)H_sA_\theta g_\Psi=0 .
 \label{eq:pair-invariance-amplitude}
\end{equation}
The defect weight from $b$ is $\sin^2(4\theta)\sin^2((j-s)t/(2\hbar))$. For bare transfer at $\theta=\pi/8$, the evolved state equals $d$ at $t=\pi\hbar/j$. During the first one percent of that wait, the orbit of the prepared plane already spans $\Hmin$ by Proposition~\ref{prop:semibounded-waiting}, while the weight on $d$ from $b$ reaches only $\sin^2(\pi/200)\approx 0.0247\%$. At $s=j$ the excited lines merge and $d$ remains orthogonal to the entire orbit. With the same preparation and the full codomain, the two decompositions are
\begin{center}
\begin{tabular}{@{}llll@{}}\toprule
Hamiltonian & Direct & Indirect & Invisible for all times\\\midrule
Bare transfer ($s=0$) & $\Span\{g,b\}$ & $\C d$ & $W$\\
Minimal completion ($s=j$) & $\Span\{g,b\}$ & $\{0\}$ & $\C d\oplus W$\\\bottomrule
\end{tabular}
\end{center}
The dimensions are $(2,1,3)$ and $(2,0,4)$. Choosing the response-minimal codomain gives adjoint kernels of dimensions one and zero, respectively.

\subsection{What a zero-imbalance reading measures}

Define $N_f=\sum_i f_i^*f_i$, $N_h=\sum_i h_i^*h_i$, and $Z_\theta=(N_f-N_h)/2$. The two rotated pair occupations have imbalance $+1$ and $-1$. States with one fermion in each rotated orbital have imbalance zero. The zero-outcome projection and its probability in a normalized state $\rho$ on $\mathcal F_2$ are
\begin{equation}
 I-Z_\theta^2=\Pi_d+\Pi_W,\qquad
 p_0(\rho)=\Tr(\rho \Pi_d)+\Tr(\rho \Pi_W).
 \label{eq:pair-ambient-zero}
\end{equation}
Here $\Pi_W$ projects onto the original outside-carrier space $W$. An upper bound $\varepsilon$ on its weight at readout gives $0\leq p_0(\rho)-\Tr(\rho \Pi_d)\leq\varepsilon$. Under ideal evolution by $H_s$, $W$ is reducing and its weight stays at the initial value.

For an explicit perturbation take $s=0$, $\theta=\pi/8$ and perturb the bare Hamiltonian to $H_0+\delta Q$, where $\delta$ has units of energy and $n_a=a^*a$. The number difference and a state in $W$ are
\[
 Q=n_{c_1}-n_{c_2}-n_{d_1}+n_{d_2},\qquad
 w_0=(c_1^*d_2^*-d_1^*c_2^*)\ket{\mathrm{vac}}/\sqrt2 .
\]
Then $Qm=2w_0$, $Qw_0=2m$ and $Qg=Qe=0$. The perturbation couples the two states $m,w_0$ which were degenerate at zero energy. In particular $w_0$ can now acquire population. At the ideal transfer time $t_*=\pi\hbar/j$ put $x=\pi\delta/j$. Starting from $b$, exponentiating the $m,w_0$ block gives the normalized state
\[
 \bigl[-e+\cos(2x)m-i\sin(2x)w_0\bigr]/\sqrt2 .
\]
Its weight on the original line $d$, its leakage into $W$, and the observed zero probability are
\begin{equation}
 p_d=\cos^4 x,\qquad p_W=\tfrac12\sin^2(2x),\qquad p_0=1-\sin^4x .
 \label{eq:pair-leakage-readout}
\end{equation}
\begin{center}
\begin{tabular}{@{}rrrr@{}}\toprule
$\delta/j$ & Weight on $d$ & Weight in $W$ & Zero probability\\\midrule
0 & 1.0000 & 0.0000 & 1.0000\\
0.03 & 0.9824 & 0.0176 & 0.9999\\
0.10 & 0.8181 & 0.1727 & 0.9909\\\bottomrule
\end{tabular}
\end{center}
The preparation places zero weight in $W$. At readout the error in using $p_0$ for $p_d$ is $2x^2+O(x^4)$, while the drop in zero probability is only $x^4+O(x^6)$. The leading leakage is therefore concealed when the zero count alone is used to estimate the line weight.

\section{Polar supports and matched observables}
\label{sec:polar-matched-observables}

The two nonzero outcomes of the occupation detector distinguish the two preparations. Its zero outcome records weight outside their span. Which part of this measurement follows for any preparation map, and which part uses the particular detector?

\subsection{What the polar map preserves}

Retain the densely defined closed map $A=u|A|$ of Section~\ref{sec:preparation-evolution}. Its initial support $S$ and final support $D$ satisfy
\begin{equation}
 S=(\ker A)^\perp=\overline{\Ran |A|}=\overline{\Ran A^*},\qquad
 D=(\ker A^*)^\perp=\overline{\Ran A},\qquad
 u^*u=\Pi_S,\quad uu^*=\Pi_D .
 \label{eq:polar-active-supports}
\end{equation}
Here $u$ is the bounded polar partial isometry that vanishes on $\ker A$. Its restriction $V=u|_S:S\to D$ is unitary, including when $A$ has a kernel or a nonclosed range. All probabilities below use the positive Hilbert inner products. Neither sector needs a ground vector.

\begin{proposition}[Born Rule Locus Correspondence, effect form]
\label{prop:matched-born-probabilities}
For a density operator $\rho$ on $S$ and an effect $0\leq E\leq I_S$, put $\rho_D=V\rho V^*$ and $E_D=VEV^*$ on $D$. Then
\begin{equation}
 \Tr_D(\rho_D E_D)=\Tr_S(\rho E).
 \label{eq:matched-effect-probability}
\end{equation}
Conversely $E_D$ is the unique effect on $D$ giving this equality for every input density operator. Conjugation by $V$ is a normal unital $*$-isomorphism from the bounded operators on $S$ to those on $D$.
\end{proposition}
\begin{proof}
$V^*V=I_S$ gives the trace identity by cyclicity. Testing a candidate effect on every pure input state fixes its quadratic form on $D$ because $V$ is onto. Polarization gives uniqueness. Conjugation preserves adjoints, products and increasing bounded suprema, and conjugation by $V^*$ is its inverse.
\end{proof}

For $E$ the projection onto a unit vector $\chi$ and input a unit vector $\psi$, the equality reads
\[
 |\langle\chi,\psi\rangle|^2=|\langle V\chi,V\psi\rangle|^2 .
\]
Thus $[\psi]\mapsto[V\psi]$ is a bijection of active rays preserving transition probabilities and their Fubini--Study distance $\arccos|\langle\chi,\psi\rangle|$. In particular a squared amplitude on either side refers to a specified outcome. The theorem supplies a matching outcome for each input effect. Determining which laboratory or observer procedures implement those effects requires further data.

The distinction between $A$ and $u$ matters even in two dimensions. For bounded $A$ and a state with $\Tr(A\rho A^*)>0$, normalizing the raw output gives
\[
 \frac{A\rho A^*}{\Tr(A\rho A^*)}
 =u\,\frac{|A|\rho|A|}{\Tr(|A|\rho|A|)}\,u^* .
\]
The positive factor filters the input before transport. For $A=\operatorname{diag}(1/2,1)$ the polar map is the identity. Starting with the equal mixture of the two basis states, the normalized raw output gives probabilities $(1/5,4/5)$ in that basis. Polar transport alone gives $(1/2,1/2)$.

\subsection{Completing a measurement on the output}

A positive-operator-valued measure $E$ on $S$ assigns an effect to each measurable event $\Delta$ and has total effect $I_S$. Its transport $VE(\Delta)V^*$ is a complete measurement on $D$. On the full output write $F(\Delta)$ for its extension by zero on $D^\perp$. The total of these effects is $\Pi_D$. To register departure as a separate outcome $\perp$, take
\begin{equation}
 F(\{\perp\})=I-\Pi_D,\qquad
 p_\perp(\sigma)=\Tr\bigl(\sigma(I-\Pi_D)\bigr)
 \label{eq:matched-defect-outcome}
\end{equation}
for any output density operator $\sigma$. The enlarged measurement is normalized and inherits countable additivity by unitary transport and direct sum. If the input measurement is projection-valued, so is this completion.

Paired preparations determine an ambient effect only through its compression to $D$. For example, an apparatus could assign all of $D^\perp$ to one existing outcome of a finite measurement. It would agree with the separate-outcome detector on every paired preparation. The same freedom concerns which observables are available. A realization may select a unital $C^*$-subalgebra $\mathcal M\subseteq\mathcal B(S)$ of admissible input observables. Its matched representation consists of the operators $X\oplus VXV^*$ on $S\oplus D$, with $X\in\mathcal M$. This selection specifies which questions are paired; the allowed generators on either sector require their own specification.

\subsection{The same detector after waiting}

An input state prepared by $V$ and evolved by $U(t)$ reaches the subspace $D_t=U(t)D$ through the isometry $V_t=U(t)V$. For a fixed detector with outcome effects $F_k$ on the output, its probabilities on any input state $\rho$ are given by the input measurement
\[
 M_k(t)=V_t^*F_kV_t,\qquad
 \Tr(\rho M_k(t))=\Tr(V_t\rho V_t^*F_k),\qquad \sum_k M_k(t)=I_S .
\]
On $D_t$ the compression of $F_k$ is exactly the transported effect $V_tM_k(t)V_t^*$. Thus the Born correspondence continues to hold on the evolved support. The input question answered by a fixed detector can change during the wait.

In the pair model let $Z_{\mathrm{in}}=\operatorname{diag}(1,-1)$ in the input pair-occupation basis. On $\mathcal F_2$ the measured $Z_\theta$ equals $A_\theta Z_{\mathrm{in}}A_\theta^*$. Indeed its $+1$ and $-1$ eigenvectors are the two rotated pair occupations and it vanishes on their orthogonal complement. Its spectral measurement therefore realizes precisely the separate defect outcome in Eq.~\eqref{eq:matched-defect-outcome}.

For this detector $V_t=\exp(-itH_s/\hbar)A_\theta$ and $F_0=I-\Pi_D$. Under $H_s$ the input ground stays on the ray of $g$, and the input excitation evolves from $b$ with amplitude $u_\theta v_\theta(1-e^{-i(j-s)t/\hbar})$ on $d$. Neither preparation reaches $W$. Hence the zero outcome pulls back to
\begin{equation}
 M_0(t)=\sin^2(4\theta)\sin^2\!\left(\frac{(j-s)t}{2\hbar}\right)\Pi_{e_\Psi}.
 \label{eq:pair-detector-pullback}
\end{equation}
At the minimal completion $s=j$ this effect vanishes. For bare transfer at $\theta=\pi/8$ and $t_*=\pi\hbar/j$, the ground goes to $-g$ and $b$ goes to $d$. The two rotated pair occupations have equal squared overlap with $g$, so all three input effects are
\begin{equation}
 M_{+1}(t_*)=M_{-1}(t_*)=\tfrac12\Pi_{g_\Psi},\qquad M_0(t_*)=\Pi_{e_\Psi}.
 \label{eq:pair-transfer-povm}
\end{equation}
The detector initially reads the two occupations. After this wait the probability of zero equals the initial excitation population. The other two probabilities sum to the ground population.

For example, prepare the incoherent mixture $\rho=\operatorname{diag}(0.36,0.64)$ in the input occupation basis. It has equal populations on $g_\Psi$ and $e_\Psi$. The probabilities in order $+1,-1,0$ are
\begin{equation}
 (0.36,0.64,0)\quad\text{at preparation},\qquad
 (0.25,0.25,0.50)\quad\text{at }t_* .
 \label{eq:pair-matched-probabilities}
\end{equation}
The input evolution at $t_*$ is $-I$ and preserves the original distribution. The coherent input with real amplitudes $(0.6,0.8)$ gives the same initial occupation counts. Its excitation weight is $|(0.6-0.8)/\sqrt2|^2=0.02$, so the same wait gives $(0.49,0.49,0.02)$.

Equation~\eqref{eq:matched-effect-probability} applies whenever the states are paired. For example, if $U(t)D=D$ for all real $t$, then $V^*U(t)V$ is the induced input unitary group on $S$. For a \emph{specified} input group $U_\Psi(t)$ on $S$, exact vector pairing requires $U(t)V=VU_\Psi(t)$ (regarding $D$ as a subspace of the output). Matching every input state and every effect, with zero defect probability, allows a common phase at each time. Indeed the vanishing defect probability puts every evolved preparation in $D$, and agreement of all effects gives identical pure-state projections there. The two isometries thus agree on each ray; linearity applied to sums makes the phase independent of the ray. A particular input or a restricted set of measurements can retain its probabilities under weaker conditions.

\section{Neutral graphs and the observability of relative phase}
\label{sec:neutral-graphs-phase}

The polar map pairs a vector $\psi$ with $V\psi$ and preserves the probability of each matched outcome. On the direct sum one can also ask how the two vectors are combined. Does changing their relative phase change any probability? The answer depends on which measurements the realization supplies.

\subsection{The graph selected by the polar map}

On the active carrier $\mathcal K_A=S\oplus D$, use the Hilbert inner product for probabilities. The sector symmetry $J(x\oplus y)=x\oplus(-y)$ supplies a second form, $[z,w]=\langle z,Jw\rangle$. For $\varphi\in\R/(2\pi\mathbb Z)$ write
\[
 G_\varphi\psi=2^{-1/2}(\psi\oplus e^{i\varphi}V\psi),
 \qquad \mathcal N_\varphi=\Ran G_\varphi .
\]
We use $\varphi$ for the relative phase and retain $\theta$ for the physical orbital rotation. The rays of $\mathcal N_0$ form the aligned Born Rule Locus. Since $V$ is unitary between supports, $G_\varphi^*G_\varphi=I_S$ and $[G_\varphi\psi,G_\varphi\chi]=0$ for all $\psi,\chi$. Thus graph vectors have their full Hilbert norm, while the Krein form vanishes on every pair from the same graph. In particular $[\psi]\mapsto[G_0\psi]$ gives the same projective transition probabilities as $[\psi]\mapsto[V\psi]$.

Each graph equals its Krein-orthogonal complement \emph{inside $\mathcal K_A$}, the hypermaximal neutral case of Krein geometry \cite{Bognar1974}. Indeed, $x\oplus y$ is Krein-orthogonal to the graph precisely when $x=e^{-i\varphi}V^*y$, or $y=e^{i\varphi}Vx$. Its Hilbert-orthogonal complement there is $\mathcal N_{\varphi+\pi}$. Equal component norms alone define the larger neutral cone; graph membership also requires the displayed vector relation.

The polar reflection $\mathcal R$ and the polar complex structure $\mathcal I$ on the active carrier are
\[
 \mathcal R=\begin{pmatrix}0&V^*\\ V&0\end{pmatrix},\qquad
 \mathcal I=J\mathcal R=\begin{pmatrix}0&V^*\\-V&0\end{pmatrix}.
\]
Here $\mathcal R^*=\mathcal R$, $\mathcal R^2=I$, $\mathcal I^*=-\mathcal I$ and $\mathcal I^2=-I$. Multiplying $G_\varphi$ by its Hilbert adjoint yields the graph projection
\begin{equation}
 \Pi_\varphi=G_\varphi G_\varphi^*
  =\tfrac12\!\left(I+\cos\varphi\,\mathcal R+\sin\varphi\,(-i\mathcal I)\right).
 \label{eq:polar-graph-projector}
\end{equation}
In particular $\mathcal N_0$ and $\mathcal N_\pi$ are the two eigenspaces of $\mathcal R$. On the full direct sum the zero extensions of these graph projections use $\Pi_S\oplus\Pi_D$ in place of $I$ and the zero extensions of the two polar operators.

\subsection{What a measurement can distinguish}

To vary the sector populations as well, fix a unit $\psi\in S$ and write
\[
 z_{p,\varphi}=\sqrt p\,\psi\oplus e^{i\varphi}\sqrt{1-p}\,V\psi,
 \qquad 0\leq p\leq1.
\]
The graphs use $p=1/2$. Holding $\psi$ fixed is essential: it keeps the state content the same on both sides while the sector weights vary.

\begin{proposition}[Relative phase and sector populations]
\label{prop:phase-readout}
Let a bounded Hilbert-self-adjoint operator on $S\oplus D$ have blocks
$T=\left(\begin{smallmatrix}B&C\\ C^*&L\end{smallmatrix}\right)$, with $C:D\to S$.
For fixed $0<p<1$ and $\psi$, its mean varies with $\varphi$ exactly when $\langle\psi,CV\psi\rangle\ne0$. It varies for some matched input exactly when $C\ne0$. Its mean is independent of \emph{both} $p$ and $\varphi$ for every matched input exactly when $T=B\oplus VBV^*$.
\end{proposition}
\begin{proof}
With the inner product linear in its second argument, the mean is
\begin{align*}
 \langle z_{p,\varphi},Tz_{p,\varphi}\rangle
 &=p\langle\psi,B\psi\rangle+(1-p)\langle V\psi,LV\psi\rangle\\
 &\quad+2\sqrt{p(1-p)}\,\operatorname{Re}\!\left(e^{i\varphi}\langle\psi,CV\psi\rangle\right).
\end{align*}
The last term is a nonconstant sinusoid whenever its coefficient is nonzero. If the coefficient vanishes for every $\psi$, complex polarization gives $CV=0$, hence $C=0$ since $V$ is onto $D$. With $C=0$, comparing $p=0$ and $p=1$ for every $\psi$ gives $B=V^*LV$. Substitution proves the converse.
\end{proof}

For a detector, apply the proposition to each outcome effect. For an ambient effect, first compress it to $S\oplus D$: a block coupling only to the defect contributes zero on these preparations. For a measurement of a self-adjoint observable, the full distribution is obtained by applying this test to its spectral projections. A single mean can be constant even when individual outcome probabilities vary. A sector projection illustrates the distinction between phase and population: the input projection has probability $p$ at every phase.

\subsection{Paired observables and the sector coordinate}

Write $\ket{\Psi},\ket{\Phi}$ for the basis of the sector coordinate $\C^2$, in input--output order. The unitary identification
\[
 x\oplus y\ \longmapsto\ \ket{\Psi}\otimes x+\ket{\Phi}\otimes V^*y
 \quad\text{from }S\oplus D\text{ to }\C^2\otimes S
\]
takes $G_\varphi\psi$ to $2^{-1/2}(\ket{\Psi}+e^{i\varphi}\ket{\Phi})\otimes\psi$.
It sends $J,\mathcal R,-i\mathcal I$ to $\sigma_z,\sigma_x,\sigma_y$ on the sector coordinate, each tensored with $I_S$. Their means on the graph state are therefore $0,\cos\varphi,\sin\varphi$, respectively. The latter two means together determine the phase modulo $2\pi$ when both measurements are available on repeated preparations.

For $X\in\mathcal B(S)$ write $\iota(X)=X\oplus VXV^*$. In the same coordinates $\iota(X)=I_2\otimes X$, so $\iota(X)G_\varphi=G_\varphi X$. This identifies the matrix elements on every graph. A density operator $\rho$ on $S$ gives the same mean $\Tr(\rho X)$ in $\rho$, $V\rho V^*$ and $G_\varphi\rho G_\varphi^*$ for the three corresponding operators. Taking $X$ to be an effect gives the full input probability even when each sector has weight one half. Taking the spectral projections of a bounded self-adjoint input gives the equality for every measurable outcome set; these projections are defined in $\mathcal B(S)$ even when the selected algebra $\mathcal M$ is smaller.

More generally the paired statistics on any state of $\C^2\otimes S$ depend only on its reduced state on $S$. An operator commutes with $I_2\otimes\mathcal M$ precisely when each of its four blocks commutes with $\mathcal M$. For the maximal choice $\mathcal M=\mathcal B(S)$ the blocks are scalars, so the commutant is $\mathcal B(\C^2)\otimes I_S$. The sector projections and phase operators belong to that commutant. Which, if any, can be measured? The polar map fixes their mathematical definition. A phase measurement requires a physically available cross-sector readout. The occupation detector of Section~\ref{sec:polar-matched-observables} acts within the output sector.

If every physically admissible measurement effect on $\mathcal K_A$ commutes with $J$, the two sectors obey a \emph{superselection rule} in the sense of Wick, Wightman and Wigner \cite{WickWightmanWigner1952}. Every such measurement gives the same probabilities for a coherent superposition across sectors and the mixture obtained by discarding its cross-sector coherences. An admissible effect whose two diagonal blocks differ after identification by $V$ can still distinguish sector populations. The paired algebra satisfies the commutation condition by construction.

\subsection{A polar rotation before a population measurement}

Suppose an apparatus can distinguish the two sectors. Its signed population observable is $J$, with input-sector effect $P_\Psi=(I+J)/2$ on $\mathcal K_A$. Every normalized graph state $G_\varphi\psi$ gives mean imbalance zero and input-sector probability $1/2$. Now suppose the apparatus can also apply the polar rotation $\mathcal U_\alpha=\exp(\alpha\mathcal I)$ before counting. Since $\mathcal I^2=-I$ on $\mathcal K_A$, this Hilbert-unitary rotation equals $I\cos\alpha+\mathcal I\sin\alpha$ there. On the full carrier the exponential fixes the inactive complement pointwise. Using $J\mathcal I=\mathcal R$ and $\mathcal I J=-\mathcal R$ gives
\begin{equation}
 \mathcal U_\alpha^*J\mathcal U_\alpha
 =J\cos(2\alpha)+\mathcal R\sin(2\alpha).
 \label{eq:quarter-phase-readout}
\end{equation}
At $\alpha=\pi/4$ the effective observable is $\mathcal R$. Its mean is $\cos\varphi$, and the probability of the $\Psi$ outcome is $(1+\cos\varphi)/2$. Thus the aligned graph ($\varphi=0$) goes entirely into the input sector, and the opposite graph ($\varphi=\pi$) entirely into the output sector. These are distinct from the action of the quarter-turn $\mathcal U_{\pi/2}=\mathcal I$, which exchanges the two sectors and sends $J$ to $-J$ under conjugation. The doubled angle in Eq.~\eqref{eq:quarter-phase-readout} accounts for the difference. Both the pulse and a detector that distinguishes the sectors are additional operational assumptions. A paired readout $I_2\otimes X$ commutes with the pulse and gives the same statistics as before it.

This example gives an exact test for superselection on the active carrier. Let $F$ range over the outcome effects of the available detectors. For a protocol $\Lambda$ preceding the detector, define its adjoint on effects by $\Tr(\Lambda(\rho)F)=\Tr(\rho\,\Lambda^*(F))$. Include the empty protocol and all allowed sequences; for a measurement branch use its trace-nonincreasing operation, so the trace gives the joint outcome probability. By Proposition~\ref{prop:phase-readout}, relative phase is invisible throughout these experiments exactly when
\begin{equation}
 [J,\Lambda^*(F)]=0\qquad\text{for every available protocol and outcome.}
 \label{eq:operational-phase-blindness}
\end{equation}
For a unitary control $U$, the effective effect is $U^*FU$. In particular the $\pi/4$ pulse followed by the sector count gives the effect $(I+\mathcal R)/2$, which violates Eq.~\eqref{eq:operational-phase-blindness}.

The two adjoints clarify what this condition requires. The Krein adjoint of a bounded operator is $T^\sharp=JT^*J$. A Hilbert-self-adjoint effect is also Krein-self-adjoint precisely when it commutes with $J$. The structural observables $\mathcal R$ and $-i\mathcal I$ are Hilbert-self-adjoint and Krein-skew-adjoint, whereas the rotation generator satisfies $\mathcal I^*=-\mathcal I$ and $\mathcal I^\sharp=\mathcal I$. For a Hilbert-unitary control, preservation of the Krein form means $U^*JU=J$, equivalently $UJ=JU$. Such a control takes every sector-preserving effect to another under pullback.

There is a weaker possibility for controls. The discrete swap $\mathcal U_{\pi/2}$ obeys $U^*JU=-J$ and exchanges the diagonal blocks of every sector-preserving effect in the factorized coordinates. A set of controls consisting of sector-preserving unitaries and this swap can therefore move population between sectors while satisfying Eq.~\eqref{eq:operational-phase-blindness} for every diagonal detector. Access to the intermediate $\pi/4$ pulse changes that conclusion when the sector count is available.

More precisely, if all block-diagonal effects are available and both sectors are nonzero, a unitary control preserves phase blindness for all those detectors exactly when $U^*JU=\pm J$. Indeed $[J,U^*FU]=0$ for every diagonal effect $F$ says that $UJU^*$ commutes with the full block-diagonal algebra. Its commutant consists of the operators scalar on each sector. Since $UJU^*$ has both eigenvalues $+1,-1$, it equals $J$ or $-J$; either choice suffices. Along a strongly continuous path from the identity for which every intermediate control obeys this condition, the sign stays positive. Which controls and detectors a realization supplies is therefore the physical issue. The two adjoint operations alone leave it open.

\subsection{When evolution preserves, rotates, or leaves a graph}

Consider separate sector evolutions $U_\Psi(t)$ on $S$ and $U(t)$ on $\mathcal H_\Phi$. With a fixed phase convention for $V$, suppose that at a given time
\begin{equation}
 U(t)V=e^{i\gamma(t)}VU_\Psi(t).
 \qquad\text{Then}\qquad
 (U_\Psi(t)\oplus U(t))G_\varphi=G_{\varphi+\gamma(t)}U_\Psi(t).
 \label{eq:graph-evolution-phase}
\end{equation}
This follows by evolving the two components. The scalar phase invisible in the separate state probabilities becomes a relative phase when the components are placed in direct sum. Shifting both sector energies by the same scalar would multiply the entire vector by one phase and leave the graph fixed.

The pair model gives all three possibilities with the preparation $V=A_\theta$. At the minimal completion $s=j$, the restricted input energies on $g_\Psi,e_\Psi$ are $-2j,0$ and the corresponding output energies on $g,b$ are also $-2j,0$. Thus $U(t)V=VU_\Psi(t)$ and every phase graph is preserved. Raising only the input generator by $\delta I$ multiplies the input evolution by $e^{-i\delta t/\hbar}$, giving $\gamma(t)=+\delta t/\hbar$ in Eq.~\eqref{eq:graph-evolution-phase}. For example, a shift $\delta=0.2j$ rotates the label by $0.2jt/\hbar$.

For bare transfer take $s=0$, $\theta=\pi/8$ and the input excitation $e_\Psi$. At $t_*=\pi\hbar/j$ the input evolution sends $e_\Psi$ to $-e_\Psi$ while the output sends $b$ to $d\perp D$. The total state and its squared projection onto any \emph{original} graph are
\begin{equation}
 z_* = 2^{-1/2}\bigl(-e_\Psi\oplus e^{i\varphi}d\bigr),\qquad
 \|\Pi_\eta z_*\|^2=\tfrac14\quad\text{for every }\eta .
 \label{eq:pair-graph-departure}
\end{equation}
Indeed $G_\eta^*z_*=-e_\Psi/2$. The vector has unit Hilbert norm and zero Krein square. Each sector still carries half the total norm squared, and all of the output component is on $d$. That component lies outside $D$, so varying the phase cannot put it on a graph of the original $V$. The input ground, by contrast, stays paired in this example. If the correspondence itself is evolved, the new isometry $U(t)VU_\Psi(t)^*$ has a graph containing the evolved state at its original phase label. Departure in Eq.~\eqref{eq:pair-graph-departure} is relative to the fixed preparation.

\begingroup\small
\begin{thebibliography}{9}
\bibitem{Kalman1963} R.~E. Kalman, Mathematical description of linear dynamical systems, \emph{J. SIAM Control} \textbf{1} (1963), 152--192. \href{https://doi.org/10.1137/0301010}{doi:10.1137/0301010}.
\bibitem{Hegerfeldt1998} G.~C. Hegerfeldt, Causality, particle localization and positivity of the energy, in \emph{Irreversibility and Causality}, Lecture Notes in Physics \textbf{504} (1998), 238--245. \href{https://doi.org/10.1007/BFb0106784}{doi:10.1007/BFb0106784}.
\bibitem{Gilbert1963} E.~G. Gilbert, Controllability and observability in multivariable control systems, \emph{J. SIAM Control} \textbf{1} (1963), 128--151. \href{https://doi.org/10.1137/0301009}{doi:10.1137/0301009}.
\bibitem{Singh2026PairTransfer} T.~P. Singh, \emph{Interacting Generalised Trace Dynamics: Cohesive Emergence, Jordan Spectra and the Problem of Fundamental Constants}, P10, 31 October 2026, Sec.~3.4, summarizing P8. \href{https://doi.org/10.5281/zenodo.23853851}{doi:10.5281/zenodo.23853851}.
\bibitem{Bognar1974} J.~Bogn\'ar, \emph{Indefinite Inner Product Spaces}, Springer, 1974. \href{https://doi.org/10.1007/978-3-642-65567-8}{doi:10.1007/978-3-642-65567-8}.
\bibitem{WickWightmanWigner1952} G.~C. Wick, A.~S. Wightman and E.~P. Wigner, The intrinsic parity of elementary particles, \emph{Phys. Rev.} \textbf{88} (1952), 101--105. \href{https://doi.org/10.1103/PhysRev.88.101}{doi:10.1103/PhysRev.88.101}.
\end{thebibliography}
\endgroup
\end{document}
